\section{Measurement}
\label{sec:methodology}

\subsection{Corpus}
\label{sec:corpus}

We instrument Claude Code, Anthropic's AI coding assistant, across
a single power user's sessions over approximately four months
(November 2025 through March 2026). The corpus comprises 857
sessions across 15 software projects, drawn from two sources:
a local WSL environment and an Ubuntu VM (migrated from a prior
machine).

Sessions are classified by type:
\begin{itemize}[nosep]
  \item \textbf{Main sessions} (59): Human-facing conversations
    where the user interacts directly with the assistant.
  \item \textbf{Subagent sessions} (567): Delegated tasks spawned
    by the main session to handle subtasks in parallel.
  \item \textbf{Compact sessions} (154): Context compaction
    continuations triggered when the main session approaches
    its context limit (the client automatically summarizes
    conversation history to stay within the context window).
  \item \textbf{Prompt suggestion sessions} (21): Short sessions
    for generating prompt suggestions.
\end{itemize}

The corpus represents 54{,}170 API calls and 4.45~billion effective
input tokens (input tokens + cache creation + cache read tokens).
Token usage is extracted from the API response metadata embedded in
session transcripts.

Because the live corpora continue to evolve, the paper reports a
frozen cohort (857 sessions) rather than a moving live count. We
provide a recount script (\texttt{tools/corpus\_counts.py}) to make
the counting boundary and dedup assumptions explicit.
Unless stated otherwise, all corpus tables and figures in this draft
use the frozen cohort with a fixed counting protocol
(\texttt{dedup\_by\_content}, \texttt{min\_size=10000}) so results are
stable across revisions.

\paragraph{Bias declaration.}
This is a single-user corpus from a power user who regularly
exhausts context windows across complex, multi-file engineering
tasks. We argue this is representative of the class where context
management matters most, and note that Pareto-distributed usage
means this class likely dominates total fleet token consumption.
Multi-user validation is identified as future work.

\subsection{Instruments}

We develop three instruments at different layers of the stack:

\paragraph{Probe} (\texttt{probe.py}): A streaming JSONL analyzer
that reads Claude Code's raw session transcripts. Classifies records
by type (user, assistant, tool\_result, progress), measures content
sizes, tracks tool usage, and computes per-session metrics including
amplification factor and tool overhead ratio. Requires no API calls;
operates on existing session files.

\paragraph{Proxy} (\texttt{proxy.py}): A transparent HTTP proxy
deployed between Claude Code and the Anthropic API. Intercepts
every request, capturing system prompts, tool definitions, message
arrays, and token usage. Optionally applies interventions (eviction,
trimming) before forwarding. Logs all decisions to JSONL for offline
analysis. This is the ``MMU'' in our analogy.

\paragraph{Experiment runner} (\texttt{pichay}): An experiment
framework that orchestrates paired runs under different treatment
conditions (baseline, trimmed, compact+trim), captures all
artifacts (proxy logs, session data, git state, test results),
and computes comparative metrics. The framework is self-bootstrapping:
it was built by Claude Code running through its own proxy, with
each round's logs feeding the next round's analysis.

\subsection{Treatment Conditions}

We define three treatment conditions applied at the proxy layer:

\begin{enumerate}[nosep]
  \item \textbf{Baseline}: Proxy in observation mode. Logs all
    traffic without mutation.
  \item \textbf{Trimmed}: Tool definition stubbing and skill
    deduplication applied to each request before forwarding.
  \item \textbf{Compact+Trim}: Stale result eviction (paging)
    combined with trimming.
\end{enumerate}

All treatments use the same model (Claude Opus), same temperature,
and same task prompts. The proxy's interposition is transparent to
both the client and the API---neither is aware of the mutations.

% ====================================================================
